\chapter*{АНОТАЦІЯ}
\addcontentsline{toc}{chapter}{АНОТАЦІЯ}

Ця робота присвячена розробці платформи IoT MLOps для моніторингу промислових сенсорів у реальному часі з виявленням аномалій на основі машинного навчання.

Актуальність дослідження обумовлена потребою промислових підприємств у доступних рішеннях для моніторингу обладнання. Існуючі корпоративні платформи мають високу вартість та складність впровадження, тоді як рішення з відкритим кодом потребують значних зусиль для інтеграції компонентів.

Метою роботи є створення багатокористувацької платформи, що поєднує потокову обробку даних, машинне навчання для детекції аномалій та MLOps практики на базі технологій з відкритим кодом.

Методологія включає аналітичний метод для вивчення існуючих платформ, методи машинного навчання (Isolation Forest, Random Forest, XGBoost) для детекції аномалій, експериментальний метод для тестування продуктивності та порівняльний аналіз алгоритмів.

Основні результати: розроблено мікросервісну архітектуру на базі Docker; забезпечено потокову обробку даних з низькою затримкою через Apache Kafka та Spark; реалізовано схема-на-тенанта ізоляцію зі стисненням TimescaleDB; досягнуто високу якість детекції аномалій на наборі даних Tennessee Eastman Process.

Практична значущість полягає у створенні готового рішення для промислових підприємств, що забезпечує раннє виявлення несправностей обладнання та зниження витрат на незаплановане обслуговування.

\vspace{1cm}

\noindent
\textbf{Ключові слова:} IoT, MLOps, виявлення аномалій, промисловий моніторинг, TimescaleDB, Apache Kafka, Apache Spark, машинне навчання, мікросервісна архітектура
